% Tre prototipi CDS: i regolamenti per il pitch.
%
% Le immagini le prepara figure.py (py -3 pitch/figure.py), poi si compila due volte:
%   pdflatex regolamenti.tex
\documentclass[10pt,a4paper]{article}
\usepackage[T1]{fontenc}
\usepackage[utf8]{inputenc}
\usepackage[italian]{babel}
\usepackage{tgheros}
\renewcommand{\familydefault}{\sfdefault}
\usepackage[a4paper,left=28mm,right=17mm,top=28mm,bottom=24mm]{geometry}
\usepackage{xcolor,graphicx,tikz,enumitem,array,tabularx,colortbl,eso-pic,calc,pgffor}
\usepackage[most]{tcolorbox}
\usepackage[tracking=true,expansion=false]{microtype}
\usepackage{qrcode}
\usetikzlibrary{calc,arrows.meta}

\pdfinfo{/Title (CDS - Tre prototipi: regolamenti) /Subject (Collocamento, Variante Carte, Mercato del Lavoro)}

% ---------------------------------------------------------------- colori: quelli delle carte
\definecolor{carta}{HTML}{F3ECDD}
\definecolor{cartascura}{HTML}{E5DACA}
\definecolor{petrolio}{HTML}{17474D}
\definecolor{rosa}{HTML}{FF3D8B}
\definecolor{cMA}{HTML}{D85A24}
\definecolor{cDI}{HTML}{2878C8}
\definecolor{cCO}{HTML}{D53883}
\definecolor{cOR}{HTML}{D2A51B}
\definecolor{cCR}{HTML}{BE3D36}
\definecolor{cRI}{HTML}{75513E}
\colorlet{arancio}{cMA}
\tcbset{colupper=petrolio}

% ---------------------------------------------------------------- la pagina è una carta
% Fascia a sinistra, filo verticale, testata con l'etichetta e il numero, piede: come sulle carte.
\newif\ifpaginacarta \paginacartatrue
\newcommand{\gioco}{CDS · TRE PROTOTIPI}
\newcommand{\tessera}{}
\colorlet{accento}{rosa}
\colorlet{fascia}{cartascura}
\colorlet{filo}{petrolio}
\colorlet{sufascia}{petrolio}
\newcommand{\duecifre}[1]{\ifnum#1<10 0\fi\number#1}

\AddToShipoutPictureBG{\AtPageLowerLeft{%
  \ifpaginacarta
  \begin{tikzpicture}[x=1mm,y=1mm]
    \useasboundingbox (0,0) rectangle (210,297);
    \fill[carta] (0,0) rectangle (210,297);
    \fill[fascia] (0,0) rectangle (16,297);
    \draw[filo,line width=1.7pt] (16,0) -- (16,297);
    \ifx\tessera\empty\else
      \begin{scope}[shift={(8,279.5)},rotate=-5]
        \fill[white] (-4.7,-4.7) rectangle (4.7,4.7);
        \fill[accento] (-4.1,-4.1) rectangle (4.1,4.1);
        \node[text=white,font=\fontsize{13}{13}\selectfont\bfseries,rotate=-5,inner sep=0] at (0,0) {\tessera};
      \end{scope}
    \fi
    \node[rotate=90,text=sufascia,font=\fontsize{6.5}{7}\selectfont,inner sep=0] at (8,110) {\textls[300]{\gioco}};
    \node[text=sufascia,font=\fontsize{6}{7}\selectfont\bfseries,inner sep=0] at (8,13.4) {CDS};
    \node[anchor=base west,text=petrolio,font=\fontsize{7.5}{9}\selectfont\bfseries,inner sep=0] at (28,280.6) {\textls[260]{\gioco}};
    \node[anchor=base east,text=petrolio,font=\fontsize{8.5}{9}\selectfont,inner sep=0] at (193,280.6) {\duecifre{\value{page}}};
    \draw[petrolio,line width=1.1pt] (28,277) -- (193,277);
    \draw[petrolio,line width=.4pt] (28,15.6) -- (193,15.6);
    \node[anchor=base west,text=petrolio,font=\fontsize{5.6}{7}\selectfont\bfseries,inner sep=0] at (28,11.6) {CDS · TRE PROTOTIPI};
    \node[anchor=base east,text=petrolio,font=\fontsize{5.6}{7}\selectfont,inner sep=0] at (193,11.6) {REGOLAMENTI · OTTOBRE 2026};
  \end{tikzpicture}%
  \fi}}

\pagestyle{empty}
\color{petrolio}
\setlength{\parindent}{0pt}
\setlength{\parskip}{1.9mm}
\linespread{1.1}
\raggedbottom
\setlist[itemize]{leftmargin=4.6mm,itemsep=.7mm,topsep=.6mm,parsep=0pt,partopsep=0pt,label={\color{accento}\rule[.24ex]{1.5mm}{1.5mm}}}
\setlist[enumerate]{leftmargin=5.2mm,itemsep=.7mm,topsep=.6mm,parsep=0pt,partopsep=0pt,label={\bfseries\arabic*.}}
\renewcommand{\arraystretch}{1.5}

% ---------------------------------------------------------------- titoli, etichette, note
\newcommand{\etichetta}[1]{{\fontsize{7}{9}\selectfont\bfseries\textls[170]{\MakeUppercase{#1}}}}
\newcommand{\titolone}[1]{\par\noindent{\color{accento}\rule{14mm}{2.8pt}}\par\vspace{1.4mm}\noindent{\fontsize{24}{26}\selectfont\bfseries\MakeUppercase{#1}}\par\vspace{3mm}}
\newcommand{\titolo}[1]{\par\vspace{4.2mm}\noindent{\color{accento}\rule{11mm}{2.4pt}}\par\vspace{1mm}\noindent{\fontsize{15}{17}\selectfont\bfseries\MakeUppercase{#1}}\par\vspace{1.2mm}}
\newcommand{\sottotitolo}[1]{\par\vspace{2.2mm}\noindent{\fontsize{10.5}{12.5}\selectfont\bfseries #1}\par\vspace{-.6mm}}
\newcommand{\attacco}[1]{{\fontsize{11.5}{16}\selectfont #1\par}}

\newtcolorbox{nota}[1]{enhanced,colback=cartascura,colframe=cartascura,boxrule=0pt,arc=0pt,outer arc=0pt,
  left=4.5mm,right=3.6mm,top=2.5mm,bottom=2.5mm,boxsep=0pt,borderline west={2.6pt}{0pt}{accento},
  before skip=3.2mm,after skip=3.2mm,parbox=false,
  before upper={{\color{accento}\etichetta{#1}}\par\vspace{.3mm}}}

% I passi numerati: il quadrato col numero, come nei regolamenti dei designer.
\newcounter{passo}
\newcommand{\daccapo}{\setcounter{passo}{0}}
\newcommand{\passo}[2]{\stepcounter{passo}\par\vspace{.4mm}\noindent
  \begin{minipage}[t]{7.4mm}\vspace{0pt}\tikz\node[fill=petrolio,text=white,font=\fontsize{11}{11}\selectfont\bfseries,minimum size=6.6mm,inner sep=0]{\thepasso};\end{minipage}\hspace{2.6mm}%
  \begin{minipage}[t]{\dimexpr\linewidth-10mm\relax}\vspace{.2mm}\textbf{#1}\enspace #2\end{minipage}\par\vspace{.6mm}}

% ---------------------------------------------------------------- icone e carte
\newcommand{\ico}[1]{\raisebox{-.26em}{\includegraphics[height=1.2em]{img/ico_#1}}}
\newcommand{\terna}[1]{\foreach \s in {#1}{\ico{\s}\hspace{.1em}}}
\newcommand{\spunta}{\tikz[baseline=.1mm]\draw[line width=1.2pt,accento,line cap=round,line join=round] (0,1.1mm) -- (.9mm,0) -- (2.7mm,2.8mm);}

% Una carta in una figura: \cartaAt[stile del bordo]{x,y}{rotazione}{larghezza in mm}{immagine}
\newcommand{\cartaAt}[5][]{%
  \begin{scope}[shift={(#2)},rotate=#3]
    \pgfmathsetmacro{\cartaW}{#4/2}\pgfmathsetmacro{\cartaH}{#4*88/126}\pgfmathsetmacro{\cartaR}{#4*0.04}
    \fill[black,opacity=.22,rounded corners=\cartaR mm] (-\cartaW+.6,-\cartaH-.8) rectangle (\cartaW+.6,\cartaH-.8);
    \begin{scope}
      \clip[rounded corners=\cartaR mm] (-\cartaW,-\cartaH) rectangle (\cartaW,\cartaH);
      \node[inner sep=0,outer sep=0,rotate=#3] at (0,0) {\includegraphics[width=#4mm,height=\dimexpr#4mm*88/63\relax]{img/#5}};
    \end{scope}
    \draw[petrolio,opacity=.6,line width=.3pt,rounded corners=\cartaR mm,#1] (-\cartaW,-\cartaH) rectangle (\cartaW,\cartaH);
  \end{scope}}

\newcommand{\senzasillabe}{\hyphenpenalty=10000 \exhyphenpenalty=10000 }

% Un punto di una carta disegnata con \cartaAt, dato in punti dell'immagine (756 x 1056).
\newcommand{\sullacarta}[6]{% nome, x e y del centro, larghezza, u, v
  \coordinate (#1) at ({#2+(#5/756-.5)*#4},{#3+(.5-#6/1056)*#4*88/63});}

% Un'illustrazione quadrata col filo, come sulle carte.
\newcommand{\quadro}[2]{\tikz\node[inner sep=0,draw=petrolio,line width=1pt]{\includegraphics[width=#1]{img/#2}};}

% Il meeple e i pezzi Competenza, per gli schemi del gioco di destrezza.
\newcommand{\meeple}[4][petrolio]{% [colore]{x,y della base}{scala}{rotazione}
  \begin{scope}[shift={(#2)},rotate=#4,scale=#3]
    \fill[#1,rounded corners=#3*.42mm] (-1,7.1) -- (-4.1,5.7) -- (-4.5,5) -- (-4.1,4.2) -- (-1.8,4.6) -- (-3.4,.3) -- (-3,0) -- (-.8,0) -- (0,1.9) -- (.8,0) -- (3,0) -- (3.4,.3) -- (1.8,4.6) -- (4.1,4.2) -- (4.5,5) -- (4.1,5.7) -- (1,7.1) -- cycle;
    \fill[#1] (0,8.2) circle (1.75);
  \end{scope}}
\newcommand{\pezzo}[4][0]{% [rotazione]{x,y del centro}{colore}{lettera}
  \begin{scope}[shift={(#2)},rotate=#1]
    \fill[black,opacity=.18,rounded corners=.6mm] (-2.5+.4,-2.5-.5) rectangle (2.5+.4,2.5-.5);
    \fill[#3,rounded corners=.6mm] (-2.5,-2.5) rectangle (2.5,2.5);
    \node[text=white,font=\fontsize{8}{8}\selectfont\bfseries,inner sep=0,rotate=#1] at (0,0) {#4};
  \end{scope}}
\newcommand{\pedana}[3]{\fill[petrolio,rounded corners=.5mm] (#1,#2-1.8) rectangle (#1+#3,#2);}

% ---------------------------------------------------------------- le aperture: un dorso di carta a tutta pagina
\newcommand{\apertura}[6]{% colore, numero, titolo, in breve, tipo di gioco, disegno
  \clearpage\paginacartafalse
  \AddToShipoutPictureBG*{\AtPageLowerLeft{\begin{tikzpicture}[x=1mm,y=1mm]
    \useasboundingbox (0,0) rectangle (210,297);
    \fill[#1] (0,0) rectangle (210,297);
    \draw[white,line width=1.1pt,rounded corners=5mm] (9,9) rectangle (201,288);
    \node[text=white,font=\fontsize{11}{13}\selectfont,inner sep=0] at (105,256) {\textls[420]{PROTOTIPO #2 DI 3}};
    \node[text=white,font=\fontsize{43}{45}\selectfont\bfseries,align=center,text width=176mm,anchor=north,inner sep=0] (t) at (105,246) {#3};
    \draw[white,line width=1.5pt] ($(t.south)+(-42,-7)$) -- ($(t.south)+(42,-7)$);
    \node[text=white,font=\fontsize{10}{13}\selectfont,anchor=north,inner sep=0] (f) at ($(t.south)+(0,-11.5)$) {\textls[170]{#5}};
    \node[text=white,font=\fontsize{13}{18.5}\selectfont\itshape,align=center,text width=136mm,anchor=north,inner sep=0] at ($(f.south)+(0,-8)$) {#4};
    #6
  \end{tikzpicture}}}%
  \null\clearpage\paginacartatrue}

\begin{document}

% ================================================================ COPERTINA
\paginacartafalse
\AddToShipoutPictureBG*{\AtPageLowerLeft{\begin{tikzpicture}[x=1mm,y=1mm]
  \useasboundingbox (0,0) rectangle (210,297);
  \fill[petrolio] (0,0) rectangle (210,297);
  \draw[white,line width=1.1pt,rounded corners=5mm] (9,9) rectangle (201,288);
  \node[text=white,font=\fontsize{14}{16}\selectfont,inner sep=0] at (105,252) {\textls[620]{CDS}};
  \node[text=white,font=\fontsize{52}{54}\selectfont\bfseries,inner sep=0] at (105,227) {TRE PROTOTIPI};
  \draw[white,line width=1.5pt] (58,211) -- (152,211);
  \node[text=white,font=\fontsize{10}{13}\selectfont,inner sep=0] at (105,202.5) {\textls[170]{COLLOCAMENTO · VARIANTE CARTE · MERCATO DEL LAVORO}};
  \node[text=white,font=\fontsize{13.5}{19}\selectfont\itshape,align=center,text width=140mm,inner sep=0] at (105,184) {Tre giochi sul mercato del lavoro: far incontrare persone, competenze, formazione e opportunità.};
  \cartaAt{31,82}{17}{52}{lav_03}
  \cartaAt{179,82}{-17}{52}{lav_02}
  \cartaAt{67,95}{9}{52}{amb_27}
  \cartaAt{143,95}{-9}{52}{amb_17}
  \cartaAt{105,100}{0}{54}{lav_07}
  \node[text=white,font=\fontsize{8.5}{10}\selectfont,inner sep=0] at (105,22) {\textls[260]{REGOLAMENTI · OTTOBRE 2026}};
\end{tikzpicture}}}
\null\clearpage\paginacartatrue

% ================================================================ IL PROGETTO
\titolone{Un tema, tre giochi}

\attacco{Far incontrare persone, competenze, formazione e opportunità: è il lavoro di chi si occupa di collocamento, ed è il tema di questi tre prototipi.}

Lo raccontano con tre linguaggi diversi: un gioco di carte a turni, una caccia alla carta in tempo reale, un gioco di destrezza. Il tono è lo stesso in tutti e tre: persone un po' sopra le righe, in uffici, cantieri e magazzini che chiunque riconosce.

Questo documento raccoglie i tre regolamenti, rivisti e allineati tra loro. Ogni gioco si legge da solo.

\vspace{3mm}
\newcommand{\scheda}[7]{% colore, numero, titolo, tipo, testo, dati, in una frase
  \begin{tcolorbox}[enhanced,width=\linewidth,colback=cartascura,colframe=cartascura,boxrule=0pt,arc=0pt,outer arc=0pt,
      left=4mm,right=4mm,top=0mm,bottom=4.5mm,boxsep=0pt,borderline north={3.2pt}{0pt}{#1}]
    \begin{minipage}[t][121mm][t]{\linewidth}
    \vspace{4.5mm}%
    \tikz[baseline=-1.1mm]{\begin{scope}[rotate=-5]\fill[white] (-.42,-.42) rectangle (.42,.42);\fill[#1] (-.36,-.36) rectangle (.36,.36);\node[text=white,font=\fontsize{12}{12}\selectfont\bfseries,rotate=-5,inner sep=0] at (0,0) {#2};\end{scope}}\par\vspace{2.6mm}
    {\color{#1}\fontsize{6.6}{8}\selectfont\bfseries\textls[140]{\MakeUppercase{#4}}\par}\vspace{1.2mm}
    {\fontsize{14.5}{16}\selectfont\bfseries\raggedright\MakeUppercase{#3}\par}\vspace{1.8mm}
    {\fontsize{9}{12.6}\selectfont\raggedright #5\par\vspace{2.2mm}#6\par}
    \vfill
    {\color{#1}\rule{8mm}{1.6pt}}\par\vspace{.6mm}
    {\fontsize{9}{12.4}\selectfont\itshape\raggedright #7\par}
    \end{minipage}
  \end{tcolorbox}}
\newcommand{\dato}[2]{\textbf{#1}\enspace #2\par\vspace{.5mm}}

\noindent
\begin{minipage}[t]{.318\linewidth}\scheda{rosa}{1}{Collocamento}{Carte · a turni}
  {Costruisci percorsi di formazione, accogli i lavoratori che li cercano e trova a ognuno il lavoro che quel percorso rende possibile. Più il lavoro somiglia alle sue ambizioni, più punti fai.}
  {\dato{Giocatori}{2--4}\dato{Ritmo}{un'azione per turno, cinque collocamenti a testa}\dato{Sul tavolo}{60 Lavoratori, 60 Ambiti, 4 token}}
  {Il gioco da tavolo: mano da gestire, mercato comune, set da chiudere.}
\end{minipage}\hfill
\begin{minipage}[t]{.318\linewidth}\scheda{petrolio}{2}{Variante Carte}{Carte · tempo reale}
  {Un lavoratore al centro, il mercato sparso sul tavolo. Tutti cercano insieme, con una mano sola e una carta alla volta, la formazione e il lavoro che combaciano. Il primo che ci riesce si prende il lavoratore.}
  {\dato{Giocatori}{2--4}\dato{Ritmo}{tre round lampo}\dato{Sul tavolo}{carte Mercato a due facce, carte Lavoratore}}
  {Tre lavoratori. Una mano sola. Una carta alla volta.}
\end{minipage}\hfill
\begin{minipage}[t]{.318\linewidth}\scheda{arancio}{3}{Mercato del Lavoro}{Destrezza · tempo reale}
  {Ogni giocatore fa crescere la propria Sede: per collocare un lavoratore, il suo meeple deve toccare i tre pezzi Competenza che il lavoro richiede. Quello che posi resta. La Sede si affolla, sale, e prima o poi qualcosa cade.}
  {\dato{Giocatori}{2--4}\dato{Ritmo}{6, 8 o 12 round simultanei}\dato{Sul tavolo}{meeple, 75 pezzi Competenza, 30 Carte Lavoro, 4 pedane}}
  {Sembra facile. Finché non devi far stare tutto in piedi.}
\end{minipage}

\titolo{Le sei competenze}

Sono il linguaggio comune dei due giochi di carte: ogni Lavoratore e ogni Ambito ne mostra tre. Nel gioco di destrezza le competenze diventano pezzi da toccare.

\vspace{1mm}
\newcommand{\competenza}[3]{\begin{minipage}[c]{.158\linewidth}\centering\includegraphics[width=11.5mm]{img/ico_#1}\par\vspace{1mm}{\fontsize{8.6}{10}\selectfont\bfseries #2\par}{\fontsize{7.2}{9}\selectfont #1 · #3\par}\end{minipage}}
\noindent\competenza{MA}{Manualità}{comune}\hfill\competenza{DI}{Digitale}{comune}\hfill\competenza{CO}{Comunicazione}{comune}\hfill\competenza{OR}{Organizzazione}{non comune}\hfill\competenza{CR}{Creatività}{non comune}\hfill\competenza{RI}{Ricerca}{rara}

\clearpage

% ================================================================ IL CAST
\titolone{Il cast}

\attacco{Sessanta lavoratori e trenta ambiti professionali, ognuno con un nome e una battuta.}

È il tono che tiene insieme i tre giochi: l'ironia di chi un ufficio, un cantiere o un centralino li ha visti davvero. Nessun eroe e nessuna caricatura cattiva: colleghi.

\vspace{2mm}
\newcommand{\ritratto}[3]{\begin{minipage}[t]{.292\linewidth}\quadro{\linewidth-1pt}{#1}\par\vspace{1.2mm}{\fontsize{10.5}{12}\selectfont\bfseries\MakeUppercase{#2}\par}\vspace{.2mm}{\fontsize{8.8}{11.6}\selectfont\raggedright #3\par}\end{minipage}}
\noindent\ritratto{ill_03}{L'Archistar}{Bellissimo. Cade, ma è bellissimo.}\hfill
\ritratto{ill_02}{Il Capocantiere}{Urla istruzioni che nessuno aveva chiesto.}\hfill
\ritratto{ill_07}{L'Amico di Tutti}{Sa il nome del tuo gatto. E di quello del capo.}

\vspace{4mm}
\noindent\ritratto{ill_22}{La Project Manager}{Ti ha messo in copia. Di nuovo.}\hfill
\ritratto{ill_10}{Lo Smanettone}{L'ha smontato per capire come funziona. Funzionava.}\hfill
\ritratto{ill_27}{Il Burocrate Gentile}{Manca un timbro. Gli dispiace moltissimo.}

\vspace{-1mm}
\titolo{E trenta ambiti in cui metterli al lavoro}

\newcommand{\ambitino}[3]{\begin{minipage}[t]{.212\linewidth}\quadro{\linewidth-1pt}{#1}\par\vspace{1mm}{\fontsize{9}{10.5}\selectfont\bfseries\MakeUppercase{#2}\par}\vspace{.1mm}{\fontsize{8.2}{10.6}\selectfont\raggedright #3\par}\end{minipage}}
\vspace{1mm}
\noindent\ambitino{art_14}{Segreteria}{Le faremo sapere.}\hfill
\ambitino{art_01}{Assistenza tecnica}{Era il cavo. È sempre il cavo.}\hfill
\ambitino{art_15}{Contabilità}{Mancano tre centesimi. Nessuno va a casa.}\hfill
\ambitino{art_28}{Sviluppo web}{Sul mio computer funziona.}

% ================================================================ 1. COLLOCAMENTO
\apertura{rosa}{1}{COLLOCAMENTO}{Forma i lavoratori rispettando le loro ambizioni, trova a ognuno il lavoro giusto e chiudi il collocamento prima degli altri.}{2--4 GIOCATORI · CARTE · A TURNI}{%
  \cartaAt{62,100}{-8}{72}{amb_08}
  \cartaAt{81,103}{-3.5}{72}{amb_13}
  \cartaAt{104,104}{2}{72}{lav_17}
  \cartaAt{130,102}{7.5}{72}{amb_17}
  \node[text=white,font=\fontsize{10}{13}\selectfont\itshape,inner sep=0] at (105,27) {«Le faremo sapere.»\enspace{\upshape\fontsize{7.5}{9}\selectfont\textls[180]{SEGRETERIA}}};
}
\renewcommand{\gioco}{COLLOCAMENTO}\renewcommand{\tessera}{1}
\colorlet{accento}{rosa}\colorlet{fascia}{cartascura}\colorlet{filo}{petrolio}\colorlet{sufascia}{petrolio}

\titolone{Il gioco in breve}

\attacco{Ogni giocatore gestisce un piccolo ufficio di collocamento.}

Con le carte Ambito costruisce percorsi di \textbf{Formazione}. Quando un percorso offre ciò che un \textbf{Lavoratore} desidera, lo accoglie. Poi gli trova un \textbf{Lavoro} che quel percorso rende possibile. Formazione, Lavoratore e Lavoro formano un \textbf{set}: più il Lavoro somiglia alle ambizioni del Lavoratore, più punti vale.

La fine della partita scatta quando qualcuno chiude il quinto set. Vince chi ha più punti.

\titolo{Materiale}

\begin{itemize}
  \item \textbf{60 carte Lavoratore.} Ognuna ha tre icone \textbf{ambizione} e un potere.
  \item \textbf{60 carte Ambito:} 30 ambiti, due copie di ognuno. Ogni carta ha tre icone \textbf{competenza} e appartiene a una di cinque \textbf{categorie}: Tecnica e produzione, Servizi e relazioni, Organizzazione e gestione, Creatività e comunicazione, Ricerca e innovazione.
  \item \textbf{4 token obiettivo:} due da 5 punti e due da 3.
\end{itemize}

\vspace{2mm}
\noindent\begin{tikzpicture}[x=1mm,y=1mm]
  \useasboundingbox (0,0) rectangle (165,84);
  \cartaAt{56,39}{0}{54}{lav_17}
  \cartaAt{137,39}{0}{54}{amb_17}
  \begin{scope}[font=\fontsize{7.8}{10.4}\selectfont\senzasillabe,line width=.5pt]
    % il lavoratore: le ambizioni, il nome, il numero
    \sullacarta{la}{56}{39}{54}{0}{36}\sullacarta{lb}{56}{39}{54}{0}{446}
    \draw ($(la)+(-1,0)$) -- ($(la)+(-2.2,0)$) -- ($(lb)+(-2.2,0)$) -- ($(lb)+(-1,0)$);
    \draw ($(la)!.5!(lb)+(-2.2,0)$) -- ($(la)!.5!(lb)+(-4.4,0)$);
    \node[anchor=east,align=right,text width=22.5mm,inner sep=0] at ($(la)!.5!(lb)+(-5.4,0)$) {\textbf{Tre ambizioni.}\\Quello che la formazione deve offrirgli.};
    \sullacarta{ln}{56}{39}{54}{0}{748}
    \draw ($(ln)+(8.4,0)$) -- ($(ln)+(-4.4,0)$); \fill ($(ln)+(8.4,0)$) circle (.6);
    \node[anchor=east,align=right,text width=22.5mm,inner sep=0] at ($(ln)+(-5.4,0)$) {\textbf{Nome e battuta.}};
    \sullacarta{lp}{56}{39}{54}{712}{50}
    \draw ($(lp)+(0,1.4)$) -- ($(lp)+(0,5.4)$) -- ($(lp)+(6.6,5.4)$); \fill ($(lp)+(0,1.4)$) circle (.6);
    \node[anchor=west,inner sep=0] at ($(lp)+(7.6,5.4)$) {\textbf{Numero:} dice quale potere ha.};
    % l'ambito: le competenze, la categoria
    \sullacarta{aa}{137}{39}{54}{0}{36}\sullacarta{ab}{137}{39}{54}{0}{446}
    \draw ($(aa)+(-1,0)$) -- ($(aa)+(-2.2,0)$) -- ($(ab)+(-2.2,0)$) -- ($(ab)+(-1,0)$);
    \node[anchor=east,align=right,text width=22mm,inner sep=0] at ($(aa)!.5!(ab)+(-3.4,0)$) {\textbf{Tre competenze.}\\Offerte come Formazione, richieste come Lavoro.};
    \sullacarta{ac}{137}{39}{54}{0}{962}
    \draw ($(ac)+(4,0)$) -- ($(ac)+(-2.4,0)$); \fill ($(ac)+(4,0)$) circle (.6);
    \node[anchor=east,align=right,text width=22mm,inner sep=0] at ($(ac)+(-3.4,0)$) {\textbf{Categoria.}\\Serve per i token.};
  \end{scope}
\end{tikzpicture}

\vspace{1mm}
Le competenze sono sei: \ico{MA} Manualità, \ico{DI} Digitale e \ico{CO} Comunicazione sono comuni; \ico{OR} Organizzazione e \ico{CR} Creatività si trovano meno; \ico{RI} Ricerca è rara. La stessa icona può comparire più volte su una carta.

\begin{nota}{Una carta, due usi}
Una carta Ambito si gioca come \textbf{Formazione} oppure come \textbf{Lavoro}. Come Formazione \emph{offre} le sue tre competenze; come Lavoro le \emph{richiede}.
\end{nota}

{\fontsize{8}{10.6}\selectfont Sulle carte del prototipo la fascia dei Lavoratori riporta ancora la dicitura «attitudini»: sono le ambizioni.\par}

\clearpage
% ---------------------------------------------------------------- preparazione e turno
\titolo{Preparazione}
\daccapo
\passo{Mescolate}{separatamente il mazzo dei Lavoratori e quello degli Ambiti.}
\passo{Pescate}{3 carte Ambito a testa: è la vostra mano.}
\passo{Aprite il mercato.}{Rivelate in fila tanti Lavoratori quanti sono i giocatori. I Lavoratori nuovi entreranno sempre in coda alla fila.}
\passo{Preparate i token}{accanto ai mazzi e lasciate spazio per gli scarti degli Ambiti.}

\vspace{1.5mm}
\noindent\begin{tikzpicture}[x=1mm,y=1mm]
  \useasboundingbox (0,-3) rectangle (165,60);
  \begin{scope}[font=\fontsize{6.4}{8}\selectfont\bfseries]
    \cartaAt{12,41}{0}{17}{dorso_amb}
    \node[inner sep=0] at (12,26.4) {\textls[150]{AMBITI}};
    \draw[petrolio,line width=.5pt,dash pattern=on 1.2mm off 1mm,rounded corners=.7mm] (25.5,29.1) rectangle (42.5,52.9);
    \node[inner sep=0] at (34,26.4) {\textls[150]{SCARTI}};
    \foreach \x/\y/\v in {55/46.5/5,65/46.5/5,55/36/3,65/36/3}{\fill[petrolio] (\x,\y) circle (4.1);\draw[carta,line width=.6pt] (\x,\y) circle (3.3);\node[text=carta,font=\fontsize{9}{9}\selectfont\bfseries,inner sep=0] at (\x,\y) {\v};}
    \node[inner sep=0] at (60,26.4) {\textls[150]{TOKEN}};
    \cartaAt{90,41}{0}{17}{lav_06}
    \cartaAt{111,41}{0}{17}{lav_28}
    \cartaAt{132,41}{0}{17}{lav_54}
    \cartaAt{155,41}{0}{17}{dorso_lav}
    \node[inner sep=0] at (111,57.4) {\textls[150]{IL MERCATO DEI LAVORATORI}};
    \node[inner sep=0] at (155,26.4) {\textls[150]{LAVORATORI}};
    \draw[accento,line width=.9pt,-{Stealth[length=2.2mm]}] (144.8,26.4) -- (139.2,26.4);
    \draw[accento,line width=.9pt,-{Stealth[length=2.2mm]}] (83,26.4) -- (77,26.4);
    \node[inner sep=0,font=\fontsize{6.6}{8}\selectfont] at (111,26.4) {chi è in fila da più tempo esce per primo};
    \cartaAt{18,8}{10}{15}{amb_13}
    \cartaAt{28,9.4}{0}{15}{amb_08}
    \cartaAt{38,8}{-10}{15}{amb_17}
    \node[anchor=west,inner sep=0] at (50,9) {\textls[150]{LA TUA MANO: 3 AMBITI}};
    \node[anchor=west,inner sep=0,font=\fontsize{8.4}{11}\selectfont,text width=62mm] at (102,9) {Inizia il giocatore più giovane, poi si procede in senso orario. Il limite di mano è \textbf{6 carte}.};
  \end{scope}
\end{tikzpicture}

\titolo{Il turno}

Nel tuo turno fai \textbf{una sola} di queste cinque azioni. I poteri dei Lavoratori possono concederne altre.
\daccapo
\passo{Aprire una formazione.}{Gioca una carta Ambito davanti a te: diventa una Formazione e apre una nuova pila. Puoi avere quante pile vuoi.}
\passo{Migliorare una formazione.}{Gioca una carta Ambito sotto una tua Formazione: le sue tre competenze si sommano a quelle della pila. Puoi farlo anche dopo aver inserito un Lavoratore, finché il set non è chiuso.}
\passo{Inserire un lavoratore.}{Prendi un Lavoratore dal mercato e mettilo sopra una tua pila che non ne ha ancora uno. Le competenze della pila devono comprendere \textbf{tutte e tre le sue ambizioni}, contate una per una: due icone Digitale ne chiedono due. Il mercato si ripristina subito, in coda. Poi \textbf{attivi il potere} del Lavoratore che hai preso.}
\passo{Completare un set con un lavoro.}{Gioca dalla mano una carta Ambito, come Lavoro, sopra una tua pila che ha già un Lavoratore. Le Formazioni della pila devono comprendere \textbf{tutte e tre le icone del Lavoro}. Segna i punti, controlla i token obiettivo e metti negli scarti le Formazioni: Lavoratore e Lavoro restano davanti a te.}
\passo{Pescare.}{Pesca 2 carte Ambito. Poi il mercato scorre: il Lavoratore che è in fila da più tempo va in fondo al mazzo e ne entra uno nuovo in coda.}

Se non puoi fare nessuna azione, passi. \textbf{Alla fine del turno}, se hai più di 6 carte in mano, scarta quelle in eccesso: è l'ultima cosa che fai.

\clearpage
% ---------------------------------------------------------------- il set e i punti
\titolo{Come si costruisce un set}

Le carte si sovrappongono lasciando visibili tutte le fasce delle icone. Ogni nuova Formazione va \textbf{sotto} le precedenti; il Lavoratore va sopra le Formazioni; il Lavoro va sopra il Lavoratore. In una pila le competenze in più, o ripetute, non danno fastidio.

\vspace{2mm}
\noindent\begin{tikzpicture}[x=1mm,y=1mm]
  \useasboundingbox (0,0) rectangle (165,60);
  \cartaAt{19,26}{0}{36}{amb_08}
  \cartaAt{25.57,26}{0}{36}{amb_13}
  \cartaAt{32.14,26}{0}{36}{lav_17}
  \cartaAt{38.71,26}{0}{36}{amb_17}
  \foreach \x/\n in {4.3/1,10.87/2,17.44/3,24.01/4}{\fill[petrolio] (\x,56.4) circle (2.2);\node[text=white,font=\fontsize{7.5}{8}\selectfont\bfseries,inner sep=0] at (\x,56.4) {\n};\draw[petrolio,line width=.4pt] (\x,54.2) -- (\x,52);}
  \begin{scope}[shift={(66,0)}]
    \newcommand{\riga}[4]{% y, numeri, titolo, testo
      \node[anchor=north west,inner sep=0,text width=97mm,font=\fontsize{8.6}{12}\selectfont] at (0,#1) {{\fontsize{7}{9}\selectfont\bfseries\color{accento}\textls[160]{#2}}\\[-.2mm]\textbf{#3}\enspace #4};}
    \riga{59}{1 · 2\enspace FORMAZIONI}{Ristorazione e Logistica.}{La pila offre \terna{MA,CO,CO}\,\terna{MA,DI,OR}}
    \riga{45.5}{3\enspace LAVORATORE}{Il Tecnico Paziente}{desidera \terna{MA,DI,CO}\,e la pila le ha tutte \spunta}
    \riga{32}{4\enspace LAVORO}{Direzione}{richiede \terna{CO,MA,OR}\,e la pila le ha tutte \spunta}
    \riga{18.5}{PUNTI}{Due ambizioni combaciano col Lavoro,}{\ico{MA} e \ico{CO}: il set vale \textbf{5 + 5 = 10 punti}. Con Assistenza clienti \terna{DI,CO,MA} varrebbe 16.}
  \end{scope}
\end{tikzpicture}

Chiuso il set, le Formazioni vanno negli scarti. Davanti a te restano il Lavoratore e il suo Lavoro: un collocamento riuscito.

\titolo{Punti del collocamento}

Ogni set vale \textbf{5 punti}, più un bonus per le ambizioni del Lavoratore che combaciano con le icone della carta giocata come Lavoro. Ogni icona si abbina una volta sola.

\vspace{1mm}
\noindent{\arrayrulecolor{carta}\setlength{\arrayrulewidth}{1.2pt}\renewcommand{\arraystretch}{1.3}%
\begin{tabularx}{\linewidth}{>{\centering\arraybackslash}X|>{\centering\arraybackslash}X|>{\centering\arraybackslash}X|>{\centering\arraybackslash}X}
\rowcolor{petrolio}\textcolor{white}{\etichetta{Corrispondenze}} & \textcolor{white}{\etichetta{Punti base}} & \textcolor{white}{\etichetta{Bonus}} & \textcolor{white}{\etichetta{Totale}}\\\hline
\rowcolor{cartascura} nessuna & 5 & 0 & \textbf{5}\\\hline
\rowcolor{cartascura} 1 & 5 & 2 & \textbf{7}\\\hline
\rowcolor{cartascura} 2 & 5 & 5 & \textbf{10}\\\hline
\rowcolor{cartascura} 3 & 5 & 11 & \textbf{16}
\end{tabularx}}

\vspace{1.5mm}
Le ambizioni non devono per forza combaciare col Lavoro: un lavoro qualunque si trova sempre. È il lavoro dei sogni che fa la differenza.

\titolo{I token obiettivo}

Per gli obiettivi conta la \textbf{categoria} della carta che ha chiuso ogni set come Lavoro. Ogni giocatore può prendere al massimo un token per obiettivo; i token si sommano al punteggio finale.

\noindent\begin{minipage}[t]{.47\linewidth}
\sottotitolo{Tre set della stessa categoria}

Il primo giocatore che chiude 3 set della stessa categoria prende il token da \textbf{5 punti}. Il secondo che ci riesce prende quello da \textbf{3 punti}.
\end{minipage}\hfill
\begin{minipage}[t]{.47\linewidth}
\sottotitolo{Quattro set di categorie diverse}

Il primo giocatore che chiude 4 set di quattro categorie diverse prende il token da \textbf{5 punti}. Il secondo che ci riesce prende quello da \textbf{3 punti}.
\end{minipage}

\clearpage
% ---------------------------------------------------------------- poteri, fine della partita
\titolo{I poteri dei lavoratori}

Ogni Lavoratore ha uno di quattro poteri. Il potere si attiva \textbf{quando inserisci il Lavoratore} con l'azione 3.

\vspace{1mm}
\noindent{\arrayrulecolor{carta}\setlength{\arrayrulewidth}{1.2pt}%
\begin{tabularx}{\linewidth}{>{\centering\arraybackslash}p{8mm}|>{\raggedright\arraybackslash}p{19mm}|>{\raggedright\arraybackslash}X}
\rowcolor{petrolio}\textcolor{white}{\etichetta{N.}} & \textcolor{white}{\etichetta{Potere}} & \textcolor{white}{\etichetta{Effetto}}\\\hline
\rowcolor{cartascura}\textbf{1} & \textbf{Pesca} & Pesca 1 carta Ambito.\\\hline
\rowcolor{cartascura}\textbf{2} & \textbf{Prenota} & Prendi un Lavoratore dal mercato e tienilo da parte, scoperto: è prenotato per te. Il mercato si ripristina.\\\hline
\rowcolor{cartascura}\textbf{3} & \textbf{Forma} & Gioca subito fino a 2 carte Ambito dalla mano come Formazioni: per aprire pile nuove o per migliorarne di tue.\\\hline
\rowcolor{cartascura}\textbf{4} & \textbf{Piazza} & Metti subito un altro Lavoratore, dal mercato o tra i tuoi prenotati, su una tua pila libera, ignorando una delle sue tre ambizioni. Il suo potere non si attiva.
\end{tabularx}}

\vspace{2mm}
\begin{itemize}
  \item \textbf{Un Lavoratore prenotato} si inserisce con l'azione 3, al posto di prenderne uno dal mercato: valgono gli stessi requisiti e il mercato non cambia. Il suo potere si attiva in quel momento, non quando lo prenoti.
  \item I Lavoratori prenotati non contano nel limite di mano e può usarli solo chi li ha prenotati.
  \item Un potere che non puoi o non vuoi usare va perso.
\end{itemize}

\begin{nota}{Chi ha quale potere}
Il potere più utile sta sui Lavoratori più difficili da formare. Nel prototipo i poteri non sono ancora stampati sulle carte: si leggono dal numero del Lavoratore.\par\vspace{1.4mm}
{\fontsize{8.4}{12}\selectfont
\textbf{1 · Pesca}\enspace 08 · 10 · 12 · 13 · 15 · 18 · 23 · 24 · 33 · 34 · 35 · 49 · 55 · 57 · 58\par
\textbf{2 · Prenota}\enspace 04 · 07 · 09 · 14 · 19 · 28 · 38 · 42 · 44 · 46 · 47 · 52 · 53 · 54 · 60\par
\textbf{3 · Forma}\enspace 05 · 16 · 17 · 25 · 26 · 32 · 37 · 39 · 43 · 45 · 48 · 50 · 51 · 56 · 59\par
\textbf{4 · Piazza}\enspace 01 · 02 · 03 · 06 · 11 · 20 · 21 · 22 · 27 · 29 · 30 · 31 · 36 · 40 · 41\par}
\end{nota}

\titolo{Fine della partita}

La fine della partita scatta quando un giocatore completa il suo \textbf{quinto set}. Terminate il giro in corso, fino al giocatore alla destra di chi ha iniziato la partita; poi giocate \textbf{ancora un giro completo}. In questi turni si possono chiudere altri set e prendere i token rimasti.

\begin{nota}{Esempio}
Chi ha iniziato la partita chiude il quinto set. Gli altri terminano il giro; poi tutti, lui compreso, giocano un ultimo turno.
\end{nota}

Quando il mazzo degli Ambiti finisce, rimescolate gli scarti. Se un mazzo finisce e non si può ricostituire, la fine della partita scatta allo stesso modo.

\sottotitolo{Vittoria}

Ognuno somma i punti dei propri set e quelli dei token. Vince chi ha il punteggio più alto. In caso di parità, tra i giocatori a pari punti vince \textbf{il primo che urla «Lavoro!»}.

% ================================================================ 2. VARIANTE CARTE
\apertura{petrolio}{2}{VARIANTE CARTE}{Tre lavoratori. Una mano sola. Una carta alla volta. Trova la formazione giusta, trova il lavoro giusto e prenditi il lavoratore prima degli altri.}{2--4 GIOCATORI · CARTE · TEMPO REALE}{%
  \cartaAt{30,120}{-24}{34}{amb_25}
  \cartaAt{58,82}{16}{34}{amb_13}
  \cartaAt{84,126}{8}{34}{amb_19}
  \cartaAt{34,66}{-9}{34}{amb_22}
  \cartaAt{150,126}{-13}{34}{amb_10}
  \cartaAt{178,100}{21}{34}{amb_28}
  \cartaAt{140,70}{-27}{34}{amb_07}
  \cartaAt{172,62}{6}{34}{amb_14}
  \cartaAt{118,118}{31}{34}{amb_01}
  \cartaAt{82,64}{-18}{34}{amb_27}
  \cartaAt[rosa,opacity=1,line width=2.2pt]{107,90}{-3}{48}{lav_19}
  \node[text=white,font=\fontsize{10}{13}\selectfont\itshape,inner sep=0] at (105,24) {«È tutto qui. Da qualche parte.»\enspace{\upshape\fontsize{7.5}{9}\selectfont\textls[180]{ARCHIVIO}}};
}
\renewcommand{\gioco}{VARIANTE CARTE}\renewcommand{\tessera}{2}
\colorlet{accento}{rosa}\colorlet{fascia}{petrolio}\colorlet{filo}{rosa}\colorlet{sufascia}{carta}

\titolone{L'idea}

\attacco{Ogni lavoratore arriva con ciò che sa già fare, ma non sempre possiede tutto quello che serve per il lavoro giusto.}

A volte basta trovare la formazione adatta. A volte il lavoro perfetto è già lì, nascosto dall'altra parte di una carta. Il problema è che il mercato non è ordinato: offerte e percorsi formativi sono sparsi ovunque, e tutti cercano nello stesso momento. Dovrete riconoscere in fretta ciò che serve, scegliere quali carte trattenere e costruire il percorso corretto prima degli altri.

\begin{nota}{Idea centrale}
Cerca una carta alla volta con una sola mano. Se ti serve, mettila davanti a te. Se non ti serve, rimettila subito nel Mercato. Appena hai costruito una combinazione valida, prendi il Lavoratore.
\end{nota}

\sottotitolo{Scopo del gioco}

Colloca più lavoratori degli altri trovando, nel caos del Mercato, le Formazioni e il Lavoro che permettono di soddisfare tutte le competenze richieste.

\titolo{Componenti}

\noindent\begin{minipage}[c]{.50\linewidth}
\sottotitolo{Carte Mercato, a due facce}

\begin{itemize}
  \item \textbf{Fronte, Formazione:} mostra una competenza che un lavoratore può acquisire.
  \item \textbf{Retro, Lavoro:} mostra le competenze richieste per accedere a quel lavoro.
\end{itemize}

\sottotitolo{Carte Lavoratore}

Ogni Lavoratore mostra le competenze che \textbf{possiede già}. Si sommano a quelle delle Formazioni raccolte durante il round.
\end{minipage}\hfill
\begin{minipage}[c]{.46\linewidth}
\begin{tikzpicture}[x=1mm,y=1mm]
  \useasboundingbox (0,0) rectangle (76,56);
  \newcommand{\schemacarta}[3]{% x, etichetta, icone
    \begin{scope}[shift={(#1,4)}]
      \fill[black,opacity=.2,rounded corners=1.2mm] (.6,-.8) rectangle (30.6,41.2);
      \fill[carta,rounded corners=1.2mm] (0,0) rectangle (30,42);
      \begin{scope}\clip[rounded corners=1.2mm] (0,0) rectangle (30,42);\fill[petrolio] (0,0) rectangle (8,42);\fill[rosa] (8,0) rectangle (8.5,42);\end{scope}
      \draw[petrolio,line width=.4pt,rounded corners=1.2mm] (0,0) rectangle (30,42);
      \foreach \s [count=\i] in {#3}{\node[inner sep=0] at (4,44-\i*8.6) {\includegraphics[width=6mm]{img/ico_\s}};}
      \node[anchor=west,inner sep=0,font=\fontsize{5.2}{6}\selectfont\bfseries] at (10.5,38.6) {\textls[200]{MERCATO}};
      \draw[petrolio,line width=.5pt] (10.5,36.8) -- (28,36.8);
      \node[anchor=west,inner sep=0,font=\fontsize{7.6}{9}\selectfont\bfseries] at (10.5,14.5) {#2};
    \end{scope}}
  \schemacarta{2}{FORMAZIONE}{DI}
  \schemacarta{44}{LAVORO}{DI,CO,MA}
  \node[anchor=west,inner sep=0,font=\fontsize{6.6}{8.4}\selectfont,text width=18mm,align=left] at (12.5,12.5) {offre\\una competenza};
  \node[anchor=west,inner sep=0,font=\fontsize{6.6}{8.4}\selectfont,text width=18mm,align=left] at (54.5,12.5) {chiede queste\\competenze};
  \draw[accento,line width=1pt,{Stealth[length=2.2mm]}-{Stealth[length=2.2mm]}] (33.5,27) to[bend left=38] (42.5,27);
  \node[inner sep=0,font=\fontsize{6}{7}\selectfont\bfseries] at (38,34) {\textls[150]{GIRA}};
  \node[inner sep=0,font=\fontsize{6.4}{8}\selectfont\bfseries] at (17,51.5) {\textls[150]{FRONTE}};
  \node[inner sep=0,font=\fontsize{6.4}{8}\selectfont\bfseries] at (59,51.5) {\textls[150]{RETRO}};
\end{tikzpicture}
\end{minipage}

\titolo{Preparazione}
\daccapo
\passo{Sparpagliate il Mercato.}{Mescolate le carte Mercato e sparpagliatele al centro del tavolo: su qualunque lato, orientate in qualsiasi direzione, anche un po' sovrapposte.}
\passo{Preparate i Lavoratori.}{Mescolate le carte Lavoratore e formate un mazzo coperto a parte.}
\passo{Lasciate spazio}{davanti a ogni giocatore per le carte che deciderà di trattenere.}

\begin{nota}{Importante}
Il Mercato deve essere raggiungibile da tutti. Non serve ordinarlo: più è caotico, più rappresenta bene il cuore del gioco.
\end{nota}

\clearpage
% ---------------------------------------------------------------- il round
\titolo{Svolgimento della partita}

La partita è composta da \textbf{3 round}. In ogni round viene rivelato un solo Lavoratore.
\daccapo
\passo{Rivela il Lavoratore.}{Rivelate la prima carta Lavoratore e mettetela al centro, ben visibile a tutti. Da questo momento la ricerca comincia.}
\passo{Cerca nel Mercato.}{Tutti giocano contemporaneamente, e per tutti valgono le stesse regole:
  \begin{itemize}
    \item puoi usare \textbf{una sola mano};
    \item puoi prendere dal Mercato \textbf{una sola carta alla volta};
    \item puoi girare e rigirare la carta che hai preso per guardarne entrambi i lati;
    \item finché hai una carta in mano non puoi prenderne un'altra.
  \end{itemize}}
\passo{Tieni o rimetti.}{Dopo aver guardato una carta scegli: se vuoi tenerla, mettila davanti a te mostrando il lato che intendi usare; se non vuoi tenerla, rimettila subito nel Mercato. Una carta posata resta com'è fino alla fine del round: non si gira e non torna nel Mercato.}
\passo{Costruisci una combinazione valida.}{Per collocare il Lavoratore devi avere davanti a te \textbf{un Lavoro} e tutte le \textbf{Formazioni} che servono perché, sommate alle competenze del Lavoratore, coprano tutti i simboli richiesti da quel Lavoro.}

\begin{nota}{Regola chiave}
Competenze del Lavoratore + Formazioni raccolte = tutte le competenze richieste dal Lavoro.
\end{nota}

\vspace{1mm}
\noindent\begin{tikzpicture}[x=1mm,y=1mm]
  \useasboundingbox (0,0) rectangle (165,48);
  \newcommand{\cartina}[3]{% x, titolo, icone
    \begin{scope}[shift={(#1,8)}]
      \fill[black,opacity=.2,rounded corners=1.1mm] (.6,-.8) rectangle (26.3,35.2);
      \fill[carta,rounded corners=1.1mm] (0,0) rectangle (25.7,36);
      \begin{scope}\clip[rounded corners=1.1mm] (0,0) rectangle (25.7,36);\fill[petrolio] (0,0) rectangle (7,36);\fill[rosa] (7,0) rectangle (7.45,36);\end{scope}
      \draw[petrolio,line width=.4pt,rounded corners=1.1mm] (0,0) rectangle (25.7,36);
      \foreach \s [count=\i] in {#3}{\node[inner sep=0] at (3.5,37.8-\i*7.4) {\includegraphics[width=5.2mm]{img/ico_\s}};}
      \node[anchor=west,inner sep=0,font=\fontsize{6.2}{7}\selectfont\bfseries] at (9.6,13) {#2};
    \end{scope}}
  \cartaAt{19,26}{0}{30}{lav_07}
  \node[font=\fontsize{20}{20}\selectfont\bfseries,inner sep=0] at (45.5,26) {+};
  \cartina{53}{FORMAZIONE}{DI}
  \node[font=\fontsize{20}{20}\selectfont\bfseries,inner sep=0] at (87.5,26) {=};
  \cartina{96}{LAVORO}{DI,CO,MA}
  \node[anchor=west,inner sep=0,text width=37mm,font=\fontsize{8.8}{12.4}\selectfont,align=left] at (128,26) {L'Amico di Tutti ha già \ico{CO} e \ico{MA}. La Formazione gli dà \ico{DI}: il Lavoro è suo \spunta};
  \node[inner sep=0,font=\fontsize{6.4}{8}\selectfont\bfseries] at (19,1.6) {\textls[150]{IL LAVORATORE}};
  \node[inner sep=0,font=\fontsize{6.4}{8}\selectfont\bfseries] at (65.8,3.6) {\textls[150]{UNA FORMAZIONE}};
  \node[inner sep=0,font=\fontsize{6.4}{8}\selectfont\bfseries] at (108.8,3.6) {\textls[150]{UN LAVORO}};
\end{tikzpicture}

\begin{itemize}
  \item \textbf{Un simbolo, una fonte.} Se il Lavoro chiede due volte la stessa competenza servono due fonti diverse: due Formazioni, oppure una Formazione e un'icona del Lavoratore.
  \item Puoi tenere davanti a te più Lavori e più Formazioni del necessario: al controllo indichi il Lavoro che usi.
\end{itemize}

\titolo{Prendere il lavoratore}

Appena ritieni di avere davanti a te una combinazione completa, \textbf{prendi fisicamente la carta Lavoratore} dal centro del tavolo, con la mano libera. Tutti si fermano e la combinazione viene controllata.

\noindent\begin{minipage}[t]{.47\linewidth}
\sottotitolo{Combinazione corretta}

Se tutte le competenze richieste dal Lavoro sono coperte, conquisti il Lavoratore e lo metti davanti a te: vale un punto.
\end{minipage}\hfill
\begin{minipage}[t]{.47\linewidth}
\sottotitolo{Combinazione errata}

Se ne manca anche una sola, il Lavoratore torna al centro e la ricerca riprende. Tu sei fuori dal round: le tue carte restano dove sono.
\end{minipage}

\clearpage
% ---------------------------------------------------------------- fine del round e della partita
\titolo{Fine del round}

Quando un Lavoratore viene conquistato:
\daccapo
\passo{Le carte tornano al centro.}{Tutte le carte Mercato trattenute dai giocatori tornano nel Mercato.}
\passo{Si rimescola il caos.}{Le carte vengono sparpagliate di nuovo.}
\passo{Si riparte.}{Si rivela un nuovo Lavoratore e comincia subito il round successivo.}

\begin{nota}{Se nessuno può riuscirci}
Se tutti i giocatori sono fuori dal round, o se tutti sono d'accordo che la combinazione non si può più completare, il Lavoratore va in fondo al mazzo, le carte tornano nel Mercato e il round si rigioca con un Lavoratore nuovo.
\end{nota}

\titolo{Fine della partita}

Dopo che sono stati conquistati \textbf{3 Lavoratori}, la partita termina. Vince chi possiede più carte Lavoratore.

\sottotitolo{Parità}

In caso di parità rivelate un quarto Lavoratore: partecipano allo spareggio solo i giocatori in parità, e il primo che lo conquista vince la partita.

\vspace{7mm}
\noindent\begin{tikzpicture}[x=1mm,y=1mm]
  \useasboundingbox (0,0) rectangle (165,96);
  \cartaAt{15,70}{-19}{29}{amb_25}
  \cartaAt{41,74}{10}{29}{amb_03}
  \cartaAt{64,68}{-33}{29}{amb_19}
  \cartaAt{110,74}{14}{29}{amb_10}
  \cartaAt{134,70}{-8}{29}{amb_28}
  \cartaAt{148,70}{25}{29}{amb_22}
  \cartaAt{20,27}{27}{29}{amb_14}
  \cartaAt{47,22}{-11}{29}{amb_08}
  \cartaAt{121,23}{-25}{29}{amb_15}
  \cartaAt{147,29}{12}{29}{amb_07}
  \cartaAt{68,33}{6}{29}{amb_13}
  \cartaAt{106,43}{-14}{29}{amb_27}
  \cartaAt[rosa,opacity=1,line width=2pt]{85,50}{3}{38}{lav_26}
\end{tikzpicture}

{\fontsize{7.6}{10}\selectfont\centering Il Mercato a metà di un round. Nel prototipo le illustrazioni sono quelle di Collocamento.\par}

% ================================================================ 3. MERCATO DEL LAVORO
\apertura{arancio}{3}{MERCATO DEL LAVORO}{Gestire lavoratori, competenze e opportunità sembra facile. Finché non devi far stare tutto in piedi.}{2--4 GIOCATORI · DESTREZZA · TEMPO REALE}{%
  \begin{scope}[shift={(105,94)},rotate=-4]
    \fill[black,opacity=.25] (-54+1.2,-54-1.6) rectangle (54+1.2,54-1.6);
    \fill[white] (-54,-54) rectangle (54,54);
    \node[inner sep=0,rotate=-4] at (0,0) {\includegraphics[width=101mm]{img/ill_03}};
  \end{scope}
  \node[text=white,font=\fontsize{10}{13}\selectfont\itshape,inner sep=0] at (105,24) {«Bellissimo. Cade, ma è bellissimo.»\enspace{\upshape\fontsize{7.5}{9}\selectfont\textls[180]{L'ARCHISTAR}}};
}
\renewcommand{\gioco}{MERCATO DEL LAVORO}\renewcommand{\tessera}{3}
\colorlet{accento}{arancio}\colorlet{fascia}{arancio}\colorlet{filo}{petrolio}\colorlet{sufascia}{white}

\titolone{L'idea}

\attacco{Sulla carta è semplice: c'è un lavoro disponibile, ci sono tre competenze richieste e da qualche parte c'è una persona da collocare.}

Nella realtà ogni scelta si appoggia su quelle precedenti. Le sedi si riempiono, le competenze si intrecciano, i lavoratori si accumulano, e basta un gesto sbagliato perché qualcosa cada fuori dal sistema.

Ogni giocatore dirige una \textbf{Sede}. Per collocare un lavoratore deve far toccare al suo meeple i tre pezzi Competenza che il lavoro richiede. Bisogna essere rapidi nel trovare la combinazione giusta, precisi nel costruirla e abbastanza audaci da crescere verso l'alto. Ma ogni errore resta lì, visibile, e a fine partita presenta il conto.

\sottotitolo{Obiettivo}

Fai più punti degli altri risolvendo Carte Lavoro, facendo crescere la tua Sede in altezza e limitando i componenti che cadono fuori.

\titolo{Componenti}

\noindent\begin{minipage}[c]{.47\linewidth}\raggedright
\begin{itemize}
  \item \textbf{25 meeple Lavoratore.}
  \item \textbf{75 pezzi Competenza:} 15 per ciascuno dei 5 tipi di competenza.
  \item \textbf{30 Carte Lavoro,} ognuna con 3 simboli Competenza.
  \item \textbf{4 pedane Sede.}
\end{itemize}

\vspace{1mm}
I tre simboli di una Carta Lavoro indicano ciò che un nuovo lavoratore deve possedere per accedere a quel lavoro.
\end{minipage}\hfill
\begin{minipage}[c]{.48\linewidth}
\begin{tikzpicture}[x=1mm,y=1mm]
  \useasboundingbox (0,0) rectangle (78,40);
  \meeple{9,14}{1.5}{0}
  \node[inner sep=0,font=\fontsize{6.2}{8}\selectfont\bfseries] at (9,8) {\textls[150]{MEEPLE}};
  \foreach \c/\l [count=\i] in {cMA/A,cDI/B,cCO/C,cOR/D,cCR/E}{\pezzo{15+\i*7.2,21}{\c}{\l}}
  \node[inner sep=0,font=\fontsize{6.2}{8}\selectfont\bfseries] at (36.6,8) {\textls[150]{PEZZI: 5 TIPI}};
  \begin{scope}[shift={(60,5)}]
    \fill[black,opacity=.2,rounded corners=1mm] (.5,-.7) rectangle (17.5,30.3);
    \fill[carta,rounded corners=1mm] (0,0) rectangle (17,31);
    \draw[petrolio,line width=.5pt,rounded corners=1mm] (0,0) rectangle (17,31);
    \node[inner sep=0,font=\fontsize{4.8}{6}\selectfont\bfseries] at (8.5,27.2) {\textls[180]{LAVORO}};
    \draw[petrolio,line width=.4pt] (2.5,25.4) -- (14.5,25.4);
    \pezzo{8.5,20}{cMA}{A}\pezzo{8.5,13.5}{cDI}{B}\pezzo{8.5,7}{cCO}{C}
  \end{scope}
  \node[inner sep=0,font=\fontsize{6.2}{8}\selectfont\bfseries] at (68.5,0.6) {\textls[150]{CARTA LAVORO}};
\end{tikzpicture}
\end{minipage}

\begin{nota}{Simboli ripetuti}
Se una carta mostra due o tre volte la stessa competenza, servono due o tre pezzi distinti di quel tipo. Un singolo pezzo non può coprire più simboli della stessa carta.
\end{nota}

\titolo{Preparazione}
\daccapo
\passo{Scegliete una Sede.}{Ogni giocatore prende una pedana e la mette davanti a sé, vuota.}
\passo{Prendete i vostri lavoratori.}{Ognuno prende i suoi meeple e li tiene accanto alla pedana: \textbf{12 a testa in 2 giocatori, 8 in 3, 6 in 4}. Il meeple che avanza resta nella scatola.}
\passo{Create il mercato centrale.}{Mettete tutti i pezzi Competenza al centro del tavolo, alla portata di tutti.}
\passo{Preparate i lavori.}{Mescolate le 30 Carte Lavoro e formate un mazzo coperto al centro.}
\passo{Lasciate spazio}{attorno a ogni Sede: i componenti che cadranno resteranno dove finiscono.}

\clearpage
% ---------------------------------------------------------------- come si gioca
\begin{nota}{Regola fondamentale: quello che posi, resta}
Quando un componente viene appoggiato sulla tua Sede e rilasciato, non puoi più spostarlo. Puoi costruire accanto, sopra o contro i pezzi già presenti, ma non puoi correggere il passato.
\end{nota}

\titolo{Come si gioca}

La partita è una serie di round rapidissimi. Tutti giocano contemporaneamente.
\daccapo
\passo{Rivela un Lavoro.}{Girate la prima Carta Lavoro del mazzo e lasciatela visibile a tutti.}
\passo{Partite insieme.}{Da questo momento tutti possono costruire sulla propria Sede, \textbf{un componente alla volta}: lo prendi, lo posi, e solo allora ne prendi un altro. Per gareggiare per la carta, il primo componente che posi è \textbf{un tuo meeple nuovo}: uno solo per round.}
\passo{Crea il contatto.}{Il meeple nuovo deve toccare direttamente tre pezzi Competenza che corrispondono ai tre simboli della carta. Puoi usare pezzi presi dal centro e pezzi già presenti nella tua Sede, senza mai spostarli.}
\passo{Lascia tutto.}{Ogni componente che appoggi e rilasci resta sulla Sede, anche se un altro giocatore completa il lavoro prima di te. Quello che hai ancora in mano quando il round finisce torna da dove l'hai preso.}
\passo{Prendi la carta.}{Appena la tua configurazione è completa e non la stai più toccando, \textbf{metti la mano sulla Carta Lavoro}: il primo che la tocca ferma il round. Dopo un rapido controllo, se è tutto giusto prendi la carta e la metti davanti a te.}
\passo{Ripartite.}{Rivelate subito la Carta Lavoro successiva.}

\titolo{Che cosa significa «a contatto»}

\noindent\begin{minipage}[c]{.56\linewidth}
\begin{itemize}
  \item Il meeple deve toccare fisicamente e direttamente ciascuno dei tre pezzi richiesti.
  \item Il contatto può avvenire di lato, dall'alto o dal basso: il gioco usa tutte e tre le dimensioni.
  \item I pezzi devono essere rilasciati. Non puoi tenere in equilibrio una configurazione con le dita mentre fermi il round.
  \item Un pezzo Competenza può servire in round successivi ad altri lavoratori, se riescono a toccarlo senza spostarlo.
\end{itemize}
\end{minipage}\hfill
\begin{minipage}[c]{.40\linewidth}
\begin{tikzpicture}[x=1mm,y=1mm]
  \useasboundingbox (0,0) rectangle (66,42);
  \node[inner sep=0,font=\fontsize{6.4}{8}\selectfont\bfseries] at (33,39.5) {\textls[150]{CARTA LAVORO: A · B · C}};
  % valido: A e B ai lati, C sopra
  \pedana{2}{10}{28}
  \meeple{16,10}{1.2}{0}
  \pezzo{9.4,12.5}{cMA}{A}
  \pezzo{22.6,12.5}{cDI}{B}
  \pezzo{16,24.5}{cCO}{C}
  \node[inner sep=0,font=\fontsize{6.4}{8}\selectfont\bfseries,text=accento] at (16,4.4) {\textls[150]{VALIDO}};
  \node[inner sep=0,font=\fontsize{6.4}{8}\selectfont] at (16,.6) {tocca A, B e C};
  % non valido: C tocca B ma non il meeple
  \pedana{36}{10}{28}
  \meeple{47,10}{1.2}{0}
  \pezzo{40.4,12.5}{cMA}{A}
  \pezzo{53.6,12.5}{cDI}{B}
  \pezzo{58.7,12.5}{cCO}{C}
  \draw[accento,line width=1.3pt,line cap=round] (56.6,18.4) -- (60.8,22.6);\draw[accento,line width=1.3pt,line cap=round] (56.6,22.6) -- (60.8,18.4);
  \node[inner sep=0,font=\fontsize{6.4}{8}\selectfont\bfseries,text=accento] at (50,4.4) {\textls[150]{NON VALIDO}};
  \node[inner sep=0,font=\fontsize{6.4}{8}\selectfont] at (50,.6) {C tocca B, non il meeple};
\end{tikzpicture}
\end{minipage}

\titolo{Chiamata errata}

Se fermi il round ma la tua configurazione non soddisfa la carta, non ottieni il Lavoro. Il round riprende subito e tu non puoi più conquistare quella carta. Tutti i componenti già posati restano dove sono.

\begin{itemize}
  \item \textbf{Se il tuo meeple nuovo cade o si sposta} prima che tu abbia completato il contatto, per questa carta sei fuori: non puoi usarne un altro.
  \item \textbf{Se nessuno può riuscirci,} perché tutti sono fuori o tutti sono d'accordo che non si può, la carta si scarta e se ne rivela un'altra.
\end{itemize}

\clearpage
% ---------------------------------------------------------------- cadute, fine, punti
\titolo{Cadute e crolli}

La Sede diventa sempre più affollata. Prima o poi qualcosa cadrà. Quando succede, non si ripara nulla.

\begin{nota}{Un pezzo caduto resta dove è caduto}
Se un meeple o un pezzo Competenza lascia del tutto la tua Sede e finisce sul tavolo, non puoi raccoglierlo, spostarlo o riutilizzarlo. Resta lì fino alla fine della partita e vale \textbf{\textminus 1 punto}.
\end{nota}

\begin{itemize}
  \item Se un crollo fa cadere più componenti, ciascuno vale \textminus 1 punto.
  \item I componenti che si spostano o rotolano ma restano sostenuti dalla Sede, direttamente o indirettamente, restano in gioco nella nuova posizione.
  \item Un componente caduto non conta per l'altezza finale della Sede.
  \item I pezzi caduti possono rendere scomodo lo spazio attorno a te: fanno parte del disastro e non si toccano.
  \item Una Carta Lavoro conquistata resta tua anche se la costruzione che l'ha vinta crolla più tardi.
\end{itemize}

\titolo{Fine della partita}

La partita termina alla fine del round in cui un giocatore \textbf{posa il suo ultimo meeple}. Sono 12 round in 2 giocatori, 8 in 3, 6 in 4.

\titolo{Punteggio}

\noindent\begin{minipage}[c]{.44\linewidth}
{\arrayrulecolor{carta}\setlength{\arrayrulewidth}{1.2pt}%
\begin{tabularx}{\linewidth}{>{\raggedright\arraybackslash}X|>{\centering\arraybackslash}p{18mm}}
\rowcolor{cartascura} Ogni Carta Lavoro & \textbf{+1}\\\hline
\rowcolor{cartascura} Sede più alta & \textbf{+3}\\\hline
\rowcolor{cartascura} Seconda Sede più alta & \textbf{+2}\\\hline
\rowcolor{cartascura} Terza Sede più alta & \textbf{+1}\\\hline
\rowcolor{cartascura} Ogni componente caduto & \textbf{\textminus 1}
\end{tabularx}}
\end{minipage}\hfill
\begin{minipage}[c]{.52\linewidth}
\begin{tikzpicture}[x=1mm,y=1mm]
  \useasboundingbox (0,0) rectangle (86,52);
  % sede alta
  \pedana{2}{8}{22}
  \foreach \p/\c/\l in {{7,10.5}/cMA/A,{12.4,10.5}/cDI/B,{17.8,10.5}/cCO/C,{9.7,15.9}/cOR/D,{15.1,15.9}/cCR/E,{12.4,21.3}/cMA/A,{12.4,26.7}/cDI/B}{\pezzo{\p}{\c}{\l}}
  \meeple{12.4,29.2}{1.05}{0}
  \node[inner sep=0,font=\fontsize{11}{11}\selectfont\bfseries,text=accento] at (13,3) {+3};
  % sede media
  \pedana{30}{8}{22}
  \foreach \p/\c/\l in {{35,10.5}/cCO/C,{40.4,10.5}/cOR/D,{45.8,10.5}/cMA/A,{40.4,15.9}/cDI/B}{\pezzo{\p}{\c}{\l}}
  \meeple{40.4,18.4}{1.05}{0}
  \node[inner sep=0,font=\fontsize{11}{11}\selectfont\bfseries,text=accento] at (41,3) {+2};
  % sede bassa, con un caduto
  \pedana{58}{8}{22}
  \foreach \p/\c/\l in {{63,10.5}/cDI/B,{68.4,10.5}/cCR/E}{\pezzo{\p}{\c}{\l}}
  \meeple{74.5,8}{1.05}{0}
  \pezzo[24]{83,3.4}{cCO}{C}
  \node[inner sep=0,font=\fontsize{11}{11}\selectfont\bfseries,text=accento] at (67,3) {+1};
  \node[inner sep=0,font=\fontsize{8}{8}\selectfont\bfseries] at (83,9.6) {\textminus 1};
  % righello
  \draw[petrolio,line width=.4pt,dash pattern=on 1mm off .8mm] (0,39.4) -- (26,39.4);
  \draw[petrolio,line width=.4pt,dash pattern=on 1mm off .8mm] (28,28.6) -- (54,28.6);
  \draw[petrolio,line width=.4pt,dash pattern=on 1mm off .8mm] (56,18.2) -- (82,18.2);
  \node[inner sep=0,font=\fontsize{6.2}{8}\selectfont\bfseries] at (42,48) {\textls[150]{CONTA IL PUNTO PIÙ ALTO ANCORA SOSTENUTO}};
\end{tikzpicture}
\end{minipage}

\sottotitolo{Come si misura l'altezza}

A fine partita confrontate il punto più alto di ogni struttura ancora sostenuta dalla sua pedana, senza toccare niente. Se la differenza non è evidente, usate un righello. Una Sede vuota non prende bonus.

In caso di parità di altezza prevale chi ha meno componenti caduti, poi chi ha più Carte Lavoro. Se la parità resta, i giocatori alla pari prendono entrambi quel bonus e il piazzamento successivo viene saltato.

\sottotitolo{Vittoria}

Somma i punti delle Carte Lavoro, il bonus altezza e le penalità per i componenti caduti. Chi ha il totale più alto ha gestito meglio il proprio mercato del lavoro e vince la partita.

\clearpage
% ---------------------------------------------------------------- esempio
\titolo{Esempio di round}

\noindent\begin{minipage}[c]{.50\linewidth}\setlength{\parskip}{1.9mm}\raggedright
La Carta Lavoro richiede le competenze A, B e C.

Marta prende uno dei suoi meeple. Nella sua Sede c'è già un pezzo A in una posizione utile: appoggia il meeple contro A, poi prende dal centro B, lo posa accanto al lavoratore, e fa lo stesso con C.

Nel frattempo Luca ha posato il suo meeple e un pezzo B, ma non ha ancora completato la combinazione.

Marta lascia la costruzione senza sostenerla e mette la mano sulla Carta Lavoro. Il contatto c'è: la carta è sua.

Luca si ferma. Il meeple e il pezzo B che ha già rilasciato restano sulla sua Sede; il pezzo che aveva in mano torna al centro. Si gira subito la carta successiva.
\end{minipage}\hfill
\begin{minipage}[c]{.46\linewidth}
\begin{tikzpicture}[x=1.22mm,y=1.22mm]
  \useasboundingbox (4,0) rectangle (64,72);
  \node[inner sep=0,font=\fontsize{6.4}{8}\selectfont\bfseries] at (34,70) {\textls[150]{LA SEDE DI MARTA}};
  \pedana{6}{44}{56}
  \meeple[petrolio!50]{16,44}{1.2}{0}
  \pezzo{9.4,46.5}{cOR}{D}\pezzo{22.6,46.5}{cDI}{B}
  \pezzo{37.4,46.5}{cMA}{A}
  \meeple{44,44}{1.2}{0}
  \pezzo{50.6,46.5}{cDI}{B}
  \pezzo{44,58.5}{cCO}{C}
  \node[inner sep=0,font=\fontsize{6.4}{8}\selectfont\bfseries,text=accento] at (34,37) {\textls[150]{A · B · C: COMPLETO}};
  \node[inner sep=0,font=\fontsize{6.4}{8}\selectfont\bfseries] at (34,27) {\textls[150]{LA SEDE DI LUCA}};
  \pedana{6}{8}{56}
  \pezzo{16,10.5}{cCR}{E}
  \meeple{34,8}{1.2}{0}
  \pezzo{40.6,10.5}{cDI}{B}
  \node[inner sep=0,font=\fontsize{6.4}{8}\selectfont] at (34,1.6) {mancano A e C: meeple e B restano dove sono};
\end{tikzpicture}
\end{minipage}

\titolo{Il senso della costruzione}

All'inizio le competenze sono facili da collegare. Poi la Sede cresce. I pezzi utili possono trovarsi nel posto sbagliato, i nuovi lavoratori devono infilarsi tra quelli vecchi e costruire verso l'alto diventa sempre più rischioso. La stessa struttura che ti aiuta a riusare le competenze è anche ciò che rende ogni nuovo collocamento più difficile.

\begin{nota}{Cuore tematico}
Far crescere il mercato del lavoro è positivo, ma più il sistema cresce più diventa complesso da gestire. La destrezza rende questa complessità fisica, visibile e immediata.
\end{nota}

\titolo{Il round in sintesi}

\begin{itemize}
  \item Rivela 1 Carta Lavoro con 3 competenze.
  \item Tutti giocano insieme, un componente alla volta.
  \item Posa un tuo meeple nuovo e fagli toccare le 3 competenze richieste.
  \item Puoi sfruttare le competenze già presenti nella tua Sede.
  \item Tutto ciò che rilasci sulla Sede resta lì.
  \item Il primo che completa e tocca la carta la prende.
  \item Ciò che cade fuori dalla Sede resta dov'è e vale \textminus 1 punto.
\end{itemize}

\clearpage
% ================================================================ A CHE PUNTO SIAMO
\renewcommand{\gioco}{CDS · TRE PROTOTIPI}\renewcommand{\tessera}{}
\colorlet{accento}{rosa}\colorlet{fascia}{cartascura}\colorlet{filo}{petrolio}\colorlet{sufascia}{petrolio}

\titolone{A che punto siamo}

\attacco{Tre prototipi a tre stadi diversi. Per ognuno: a che punto è, e che cosa resta da decidere con i pezzi in mano.}

\vspace{2mm}
\newcommand{\stato}[5]{% colore, numero, titolo, fatto, da fare
  \noindent\begin{tcolorbox}[enhanced,colback=cartascura,colframe=cartascura,boxrule=0pt,arc=0pt,outer arc=0pt,left=5mm,right=5mm,top=3.4mm,bottom=3.4mm,boxsep=0pt,borderline west={3.2pt}{0pt}{#1},before skip=0mm,after skip=3.6mm]
    \tikz[baseline=-1.1mm]{\begin{scope}[rotate=-5]\fill[white] (-.42,-.42) rectangle (.42,.42);\fill[#1] (-.36,-.36) rectangle (.36,.36);\node[text=white,font=\fontsize{12}{12}\selectfont\bfseries,rotate=-5,inner sep=0] at (0,0) {#2};\end{scope}}%
    \hspace{2.4mm}{\fontsize{13.5}{15}\selectfont\bfseries\MakeUppercase{#3}}\par\vspace{2mm}
    \begin{minipage}[t]{.5\linewidth}{\color{#1}\etichetta{A che punto è}}\par\vspace{.6mm}{\fontsize{9}{12.6}\selectfont\raggedright #4\par}\end{minipage}\hfill
    \begin{minipage}[t]{.44\linewidth}{\color{#1}\etichetta{Da decidere al tavolo}}\par\vspace{.6mm}{\fontsize{9}{12.6}\selectfont\raggedright #5\par}\end{minipage}
  \end{tcolorbox}}

\stato{rosa}{1}{Collocamento}
  {Regole complete e carte illustrate: 60 Lavoratori e 30 Ambiti. Le regole di questo documento sono state giocate da bot in migliaia di partite simulate: la partita arriva sempre al quinto set, in 2, in 3 e in 4 giocatori, in circa 25 giri; ognuno chiude tra 4 e 5 set e fa circa 55 punti; chi comincia non ha vantaggio.}
  {I quattro poteri. Due di loro, Forma e Piazza, chiedono preparazione e nelle simulazioni vanno spesso a vuoto: vanno osservati con giocatori veri. Poi i poteri vanno stampati sulle carte.}

\stato{petrolio}{2}{Variante Carte}
  {Prototipo v0.1. Il regolamento è completo e pronto per i primi playtest: tre round, una regola chiave.}
  {Il mazzo Mercato a due facce è da realizzare. Con le carte in mano si decide quante ne servono, quante competenze chiede un Lavoro e quanto dura un round.}

\stato{arancio}{3}{Mercato del Lavoro}
  {Prototipo v0.1. Un modello al computer della partita, con migliaia di partite simulate, ha guidato la revisione: i meeple personali che contano i round, la carta toccata per fermare il gioco, un componente alla volta.}
  {La forma dei pezzi Competenza e la misura della pedana: decidono quanto la Sede si affolla, quanti lavoratori possono condividere lo stesso pezzo e quanto spesso qualcosa cade.}

\vspace{1mm}
\noindent\begin{minipage}[c]{.72\linewidth}
\titolo{Provare Collocamento adesso}

Le carte di Collocamento sono su un tavolo online, da aprire nel browser: un tavolo libero, senza regole imposte, per giocare a distanza con il regolamento alla mano.

\vspace{1mm}
{\fontsize{10.5}{13}\selectfont\bfseries collocamento.cortivo81.workers.dev}
\end{minipage}\hfill
\begin{minipage}[c]{.22\linewidth}\raggedleft
\color{petrolio}\qrcode[height=30mm,nolink]{https://collocamento.cortivo81.workers.dev}
\end{minipage}

\end{document}
